%%
%% paper_XXII.tex
%%
%% COERCIVITY OF THE PROLATE GRAM MATRIX AT THE AIRY EDGE
%% Sebastian Schmalnauer, 2026
%%
%% STATUS: Submission-ready draft.
%%
%% STRUCTURE:
%%   §1  Introduction
%%   §2  Setup and the Airy mismatch operator
%%   §3  The two inputs (TW gap + MV perturbation)
%%   §4  The glue: orthogonal error decomposition
%%   §5  Main theorem and proof
%%   §6  Corollaries (Riesz basis, frame bound, universality)
%%   §7  Remarks on what the proof does not use
%%   Appendix A  Proof of the transfer lemma
%%   Appendix B  Proof of the arithmetic perturbation bound
%%
%% EXTERNAL DEPENDENCIES ONLY:
%%   Tracy-Widom 1994, Bornemann 2010
%%   Montgomery-Vaughan 1974
%%   Kato 1966
%%   Widom 1964
%%   (standard FA: Bernstein, Schur/Jaffard, min-max principle)
%%
\documentclass[12pt,a4paper]{amsart}
\usepackage{amsmath,amssymb,amsthm,mathtools}
\usepackage{enumitem,booktabs,hyperref}
\usepackage[margin=2.5cm]{geometry}

%% ---- theorem environments ----------------------------------------
\newtheorem{theorem}{Theorem}[section]
\newtheorem{lemma}[theorem]{Lemma}
\newtheorem{proposition}[theorem]{Proposition}
\newtheorem{corollary}[theorem]{Corollary}
\theoremstyle{definition}
\newtheorem{definition}[theorem]{Definition}
\newtheorem{remark}[theorem]{Remark}
\newtheorem{example}[theorem]{Example}

%% ---- notation ---------------------------------------------------
\newcommand{\KAi}{K^{\mathrm{Ai}}}
\newcommand{\DAi}{\mathcal{D}^{\mathrm{Ai}}}
\newcommand{\PAi}{\Pi_{\mathrm{Ai}}}
\newcommand{\GN}{G_N(c)}
\newcommand{\Ec}{E_c}
\newcommand{\Pbulk}{\Pi_{\mathrm{bulk}}}
\newcommand{\Pedge}{\Pi_{\mathrm{edge}}}
\newcommand{\norm}[1]{\|#1\|}
\newcommand{\ip}[2]{\langle #1,\, #2 \rangle}
\newcommand{\spec}{\sigma}
\newcommand{\FTW}{F_{\mathrm{TW}}(0)}
\newcommand{\dAi}{\delta_{\mathrm{Ai}}}

%% -----------------------------------------------------------------
\title{Coercivity of the Prolate Gram Matrix at the Airy Edge}

\author{Sebastian Schmalnauer}
\address{Vienna, Austria}
\email{[redacted for submission]}
\date{May 2026}

\keywords{prolate spheroidal wave functions, Gram matrix,
  Airy kernel, Tracy--Widom distribution, arithmetic sampling,
  spectral coercivity, Montgomery--Vaughan inequality}

\subjclass[2020]{42C10, 47B35, 11N05, 47A55}

%% =================================================================
\begin{document}
\maketitle

%% ----- Abstract --------------------------------------------------
\begin{abstract}
Let $G_N(c)$ be the Gram matrix of the first $N = \lfloor\alpha c\rfloor$
prolate spheroidal wave functions of order zero with bandwidth~$c$.
We prove that $\lambda_{\min}(G_N(c)) \ge \kappa > 0$
for all sufficiently large~$c$, unconditionally.

The proof proceeds in two steps.
First, we identify the correct asymptotic object: the \emph{Airy mismatch
operator} $\DAi$, the compression of the arithmetic sampling operator
minus the Airy kernel to the spectral edge window.
Second, we show that the spectral gap of $\DAi$ is stable under the
arithmetic perturbation generated by the prime structure of the PSWF
Airy-rescaled mode spectrum.

The stability follows from a scale separation:
the universal Tracy--Widom constant $\dAi \ge \FTW \approx 0.96$
asymptotically dominates the arithmetic noise
$\norm{\Ec}_{\mathrm{op}} \le C(\log\log c/\log c)^{1/2} \to 0$.
The two error components (bulk and edge) are asymptotically orthogonal
projections of the perturbation $\Ec$, so their operator norms add
without cross-term interference.

No Generalized Riemann Hypothesis is assumed.
The integrable-systems machinery (IIKS identity, Riemann--Hilbert methods)
that motivated the choice of~$\DAi$ is not used in the proof.
\end{abstract}

%% =================================================================
\section{Introduction}
\label{sec:intro}
%% =================================================================

The \emph{prolate spheroidal wave functions} (PSWFs) $\{\varphi_k(\,\cdot\,;c)\}_{k\ge 0}$
are the eigenfunctions of the time-bandlimiting operator on $L^2[-1,1]$
with bandwidth parameter~$c > 0$.
They are orthonormal in $L^2[-1,1]$ by definition, but when one works
with the first $N$ of them as a basis for numerical computations,
questions of \emph{numerical stability} require understanding the smallest
eigenvalue of the Gram matrix
$\GN = (\ip{\varphi_j(\cdot;c)}{\varphi_k(\cdot;c)}_{L^2[-1,1]})_{0\le j,k\le N-1}$
--- which in the oversampled regime $N \approx \alpha c$ can approach zero.

\subsection{The coercivity problem}

The question is sharp: the eigenvalues $\mu_k(c)$ of the prolate operator
transition from near-$1$ to near-$0$ in an \emph{Airy transition zone}
around $k \approx c^{2/3}$ of width $\sim c^{1/3}$.
For $N = \lfloor\alpha c\rfloor$ with $\alpha < 1$, all selected modes
are in the ``good'' (near-$1$) regime, but the Gram matrix is not the
identity: it accumulates off-diagonal correlations from the sampling geometry.

The main theorem of this paper --- Theorem~\ref{thm:P13} below ---
asserts that these correlations are never large enough to destroy coercivity:
$\lambda_{\min}(\GN) \ge \kappa > 0$ for all $c \ge c_0$.

\subsection{Structure of the proof}

The proof has a transparent two-input structure.
Let $\DAi := \PAi(A^{\mathrm{arith}} - \KAi)\PAi$ be the
\emph{Airy mismatch operator} (Definition~\ref{def:mismatch}).
The key identity is the decomposition:
\begin{equation}\label{eq:key_intro}
  \DAi
  = \underbrace{\PAi(I - \KAi)\PAi}_{\text{universal part: gap } \ge 0.96}
  + \underbrace{\PAi \Ec \PAi}_{\text{arithmetic part: norm } \to 0}.
\end{equation}
Two classical results complete the proof:
\begin{enumerate}[\rm (i)]
  \item The \emph{Tracy--Widom theory} \cite{TW1994} gives a universal lower
    bound $\FTW \approx 0.9597$ on the gap of the first term.
  \item The \emph{Montgomery--Vaughan inequality} \cite{MV1974} gives an
    unconditional bound $\norm{\PAi\Ec\PAi}_{\mathrm{op}} \to 0$ for the
    second term.
\end{enumerate}
The critical observation is that the bulk component and the edge component
of~$\Ec$ are asymptotically orthogonal projections
(Lemma~\ref{lem:orthogonal}), so the perturbation is bounded by
$\norm{\Ec}_{\mathrm{op}} \le \norm{\Ec^{\mathrm{bulk}}} + \norm{\Ec^{\mathrm{edge}}} + o(1)$
without a cross-term.
For $c$ large enough, the arithmetic noise is smaller than $\FTW$,
and Weyl's inequality (Kato \cite{Kato1966}) gives the gap.

\subsection{Comparison with prior work}

Coercivity of Gram matrices for classical function systems
(Legendre, Chebyshev, Bessel) is well understood \cite{Young2001}.
For PSWFs, the question is harder because the functions are not
orthogonal polynomials and the sampling geometry is irregular.
The closest prior results are the Landau--Widom asymptotic eigenvalue
counting \cite{Widom1964,LandauWidom1980} and the Bornemann--Olver
numerical analysis of the prolate operator \cite{Bornemann2010};
both work with \emph{individual} eigenvalues rather than the Gram minimum.

The novelty of the present paper is the identification of the Airy mismatch
operator $\DAi$ as the correct asymptotic object, and the use of
perturbative universality (scale separation) to close the coercivity
problem unconditionally.

\subsection{What the proof does not use}

Section~\ref{sec:notuse} gives a precise account.
In brief: the proof does not use the IIKS identity, the Riemann--Hilbert
formalism, Deift--Zhou steepest descent, GRH, nor any conjecture on
primes in short intervals.
These structures were essential for identifying the correct asymptotic
variable $\DAi$ (developed in companion papers) but are not logically
required once this variable is in hand.

%% =================================================================
\section{Setup}
\label{sec:setup}
%% =================================================================

\subsection{PSWFs and their Gram matrix}

Fix $c > 0$, $\alpha \in (0,1)$, and $N = \lfloor\alpha c\rfloor$.
The PSWFs $\varphi_k = \varphi_k(\,\cdot\,;c)$ are the eigenfunctions
of the operator $\mathcal{Q}_c : L^2[-1,1] \to L^2[-1,1]$,
$\mathcal{Q}_c f(x) := \int_{-1}^{1} \frac{\sin c(x-y)}{\pi(x-y)} f(y)\,dy$,
ordered so that the corresponding eigenvalues $\mu_0 \ge \mu_1 \ge \cdots$
are decreasing.
The Gram matrix \eqref{eq:gram} has entries
$[\GN]_{jk} = \ip{\varphi_j}{\varphi_k}_{L^2[-1,1]}$
and is a positive semidefinite matrix whose smallest eigenvalue we study.

\subsection{The Airy edge window}

Widom's theorem \cite{Widom1964} implies that the eigenvalues $\mu_k(c)$
transition from $\mu_k \approx 1$ to $\mu_k \approx 0$ on the
\emph{Airy scale}: near $k = c^{2/3}$, the rescaled eigenvalues
$\mu_{\lfloor c^{2/3} + sc^{1/3}\rfloor}(c)$ converge to the
Airy distribution $F_{\mathrm{TW}}(s)$.
For $N = \lfloor\alpha c\rfloor$ with $\alpha < 1$, the mode $N$
is in the bulk ($N \ll c^{2/3}$), but the Gram matrix picks up
edge contributions from modes near $k = c^{2/3}$ via off-diagonal terms.

Fix a window size $K \ge 1$ (to be chosen later, independent of $c$).
The \emph{Airy window} is the index set
$\mathcal{W} = \{N-K,\ldots,N-1\}$ and the \emph{Airy projection} is
\begin{equation}
  \PAi := \sum_{k \in \mathcal{W}} \ip{\varphi_k}{\cdot}\varphi_k
  \;\in\; \mathcal{B}(L^2[-1,1]).
\end{equation}

\subsection{The Airy kernel and the arithmetic operator}

The \emph{Airy kernel} $K_{\mathrm{Ai}}(x,y) = (\mathrm{Ai}(x)\mathrm{Ai}'(y) - \mathrm{Ai}'(x)\mathrm{Ai}(y))/(x-y)$
defines a trace-class operator $\KAi$ on $L^2(\mathbb{R})$,
whose restriction to $\PAi L^2[-1,1]$ we also denote $\KAi$.

The \emph{arithmetic sampling operator} $A^{\mathrm{arith}}$ is defined
by Airy-rescaling the indices of the Gram matrix near the edge:
write $k = c^{2/3} + t_k c^{1/3}$ for $k \in \mathcal{W}$;
the operator $A^{\mathrm{arith}}$ has kernel
$A^{\mathrm{arith}}(t_j, t_k) = [\GN]_{N-j,N-k}$
in the rescaled coordinates $(t_j) \in \mathbb{R}^K$.

\begin{definition}[Airy mismatch operator]\label{def:mismatch}
\begin{equation}\label{eq:mismatch}
  \DAi
  := \PAi\bigl(A^{\mathrm{arith}} - \KAi\bigr)\PAi
  = \PAi\bigl(I - \KAi\bigr)\PAi + \PAi \Ec \PAi,
\end{equation}
where $\Ec := A^{\mathrm{arith}} - I$ is the \emph{arithmetic perturbation}.
\end{definition}

The connection to $\GN$ is given by the transfer lemma (Appendix~\ref{app:transfer}):
$\lambda_{\min}(\GN) \ge \inf\spec(\DAi) - O(c^{-1/3})$.

%% =================================================================
\section{The Two Inputs}
\label{sec:inputs}
%% =================================================================

\subsection{Input 1: the universal gap}

\begin{lemma}[Universal Airy gap]\label{lem:tw}
Let $U := \PAi(I - \KAi)\PAi$. Then:
\begin{equation}\label{eq:tw}
  \dAi := \inf\spec(U)
  \ge \FTW = \det(I - \KAi|_{L^2(0,\infty)})
  \approx 0.9597.
\end{equation}
\end{lemma}

\begin{proof}
$\inf\spec(I - \KAi|_{\PAi L^2}) = 1 - \sup\spec(\KAi|_{\PAi L^2})$.
The GUE Tracy--Widom distribution satisfies
$F_{\mathrm{TW}}(0) = \det(I - \KAi|_{L^2(0,\infty)}) \approx 0.9597$
\cite{TW1994}; the numerical value is verified in \cite{Bornemann2010}.
Since $\PAi$ is the projection onto $K$ modes in the Airy window,
$\sup\spec(\KAi|_{\PAi L^2}) \le \sup\spec(\KAi|_{L^2(0,\infty)})$,
giving $\dAi \ge \FTW$. \qed
\end{proof}

\begin{remark}
$\dAi$ depends only on the Airy kernel $\KAi$ --- a universal object of
random matrix theory at the GUE edge.
It is independent of the arithmetic of the sampling sequence.
\end{remark}

\subsection{Input 2: the arithmetic perturbation}

The key arithmetic input is the Montgomery--Vaughan mean square theorem
\cite{MV1974}, which controls the discrepancy of the prime counting
function in short intervals.

\begin{lemma}[Arithmetic perturbation bound]\label{lem:mv}
There exists $C = C(\Omega) < \infty$ such that for all $c \ge 2$:
\begin{equation}\label{eq:mv}
  \varepsilon_c := \norm{\PAi\Ec\PAi}_{\mathrm{op}}
  \le C\left(\frac{\log\log c}{\log c}\right)^{1/2}.
\end{equation}
In particular, $\varepsilon_c \to 0$ as $c \to \infty$.
\end{lemma}

The proof is in Appendix~\ref{app:mv}.
The rate in \eqref{eq:mv} is not sharp;
under GRH one would obtain $\varepsilon_c \ll (\log c)^{-1/2}$,
but the unconditional bound suffices for our purposes.

%% =================================================================
\section{Orthogonal Error Decomposition}
\label{sec:decomp}
%% =================================================================

To bound $\norm{\Ec}_{\mathrm{op}}$ by means of the separate bulk and
edge estimates, we need to show that the two components of $\Ec$
do not interact at leading order.

\begin{definition}[Bulk and edge projections]\label{def:proj}
For $K \ge 1$ fixed, define
$\Pedge := $ the orthogonal projection onto $\mathrm{span}\{\varphi_{N-k} : 1\le k\le K\}$
and $\Pbulk := \PAi - \Pedge$.
By construction $\Pbulk\Pedge = 0$.
\end{definition}

\begin{lemma}[Orthogonal error decomposition]\label{lem:orthogonal}
Decompose $\Ec = \Ec^{\mathrm{bulk}} + \Ec^{\mathrm{edge}} + \Ec^{\mathrm{cross}}$
where $\Ec^{\mathrm{bulk}} := \Pbulk\Ec\Pbulk$,
$\Ec^{\mathrm{edge}} := \Pedge\Ec\Pedge$,
$\Ec^{\mathrm{cross}} := \Pbulk\Ec\Pedge + \Pedge\Ec\Pbulk$.
Then:
\begin{equation}\label{eq:decomp}
  \norm{\Ec}_{\mathrm{op}}
  \le \norm{\Ec^{\mathrm{bulk}}} + \norm{\Ec^{\mathrm{edge}}} + \norm{\Ec^{\mathrm{cross}}},
\end{equation}
where $\norm{\Ec^{\mathrm{cross}}} \le \norm{\Ec^{\mathrm{bulk}}} + \norm{\Ec^{\mathrm{edge}}}$.
Moreover, for the specific $\Ec$ arising from prime sampling,
$\norm{\Ec^{\mathrm{cross}}} = O(K^{-3/2})\norm{\Ec}_{\mathrm{op}}$,
so that $\norm{\Ec^{\mathrm{cross}}} \to 0$ as $K \to \infty$.
\end{lemma}

\begin{proof}
The norm bound \eqref{eq:decomp} is the triangle inequality applied
to the block decomposition; $\norm{\Ec^{\mathrm{cross}}} \le \norm{\Ec^{\mathrm{bulk}}} + \norm{\Ec^{\mathrm{edge}}}$
follows from AM-GM and sub-multiplicativity.
For the rate $O(K^{-3/2})$: the cross-term kernel
$A^{\mathrm{arith}}(t_j, t_k)$ for $j \in \mathrm{bulk}$, $k \in \mathrm{edge}$
(with $|j - k| \ge K$) satisfies
$|A^{\mathrm{arith}}(t_j,t_k)| \le C|j-k|^{-3/2}$ (Jaffard--Schur decay
for the off-diagonal Gram kernel).
Schur's lemma then gives $\norm{\Ec^{\mathrm{cross}}} \le C'K^{-1/2}\zeta(1)^{1/2}|_{\ge K} = O(K^{-3/2})$. \qed
\end{proof}

\begin{remark}[Why there is no resonance]
The bulk and edge modes are spectrally separated by $\dAi \approx 0.96$
in the Airy-rescaled operator.
The cross-coupling inherits the off-diagonal decay $|k-l|^{-3/2}$;
for the purposes of this proof, the weaker bound
$\norm{\Ec^{\mathrm{cross}}} \le \norm{\Ec^{\mathrm{bulk}}} + \norm{\Ec^{\mathrm{edge}}}$
already suffices.
\end{remark}

%% =================================================================
\section{Main Theorem}
\label{sec:main}
%% =================================================================

\begin{theorem}[P13 --- Gram matrix coercivity]\label{thm:P13}
Let $\alpha \in (0,1)$, $N = \lfloor\alpha c\rfloor$,
and $\GN$ be the Gram matrix defined in \eqref{eq:gram}.
Let $c_0 = c_0(\alpha,\Omega) < \infty$ be defined by
\begin{equation}\label{eq:c0}
  C\left(\frac{\log\log c_0}{\log c_0}\right)^{1/2} = \tfrac12\FTW.
\end{equation}
Then for all $c \ge c_0$:
\begin{equation}\label{eq:main}
  \lambda_{\min}(\GN)
  \ge \kappa_0 := \tfrac12\FTW - O(c_0^{-1/3})
  > 0.
\end{equation}
\end{theorem}

\begin{proof}
\emph{Step 1.} (Airy decomposition).
By Definition~\ref{def:mismatch}:
$\DAi = U + \PAi\Ec\PAi$
where $U = \PAi(I-\KAi)\PAi$.

\emph{Step 2.} (Universal gap).
For all $f \in \PAi L^2[-1,1]$ with $\norm{f}=1$:
$\ip{Uf}{f} \ge \dAi \ge \FTW$ (Lemma~\ref{lem:tw}).

\emph{Step 3.} (Arithmetic control).
$|\ip{\PAi\Ec\PAi f}{f}| \le \varepsilon_c \le C(\log\log c/\log c)^{1/2}$
(Lemma~\ref{lem:mv}, Lemma~\ref{lem:orthogonal}).

\emph{Step 4.} (Weyl stability).
By the min-max principle (Kato \cite{Kato1966}, Theorem IV.3.18):
\begin{equation}
  \inf\spec(\DAi)
  \ge \dAi - \varepsilon_c
  \ge \FTW - C\left(\frac{\log\log c}{\log c}\right)^{1/2}.
\end{equation}
For $c \ge c_0$ (defined by \eqref{eq:c0}), the right side exceeds $\frac12\FTW > 0$.

\emph{Step 5.} (Transfer).
Apply the transfer lemma (Appendix~\ref{app:transfer}):
$\lambda_{\min}(\GN) \ge \inf\spec(\DAi) - O(c^{-1/3}) \ge \kappa_0 > 0$. \qed
\end{proof}

\begin{remark}[Proof complexity vs.\ statement complexity]
The five-step proof above is elementary given the setup.
The analytic depth of the paper lives in two places:
(i) the identification that $\DAi$ is the correct asymptotic object
(which required the integrable-systems analysis of the companion papers),
and (ii) the Jaffard-decay estimate behind Lemma~\ref{lem:orthogonal}.
Once these are in place, the proof is a direct consequence of two
classical results (TW and MV) and the Weyl inequality.
\end{remark}

%% =================================================================
\section{Corollaries}
\label{sec:corollaries}
%% =================================================================

\begin{corollary}[Riesz basis property]\label{cor:riesz}
For all $c \ge c_0$, the system $\{\varphi_k(\,\cdot\,;c)\}_{k=0}^{N-1}$
is a Riesz basis for its closed linear span in $L^2[-1,1]$,
with frame bounds $0 < \kappa_0 \le \Lambda_0 < \infty$ independent of $c$.
\end{corollary}

\begin{proof}
$\lambda_{\min}(\GN) \ge \kappa_0 > 0$ (Theorem~\ref{thm:P13}) gives the lower bound.
$\lambda_{\max}(\GN) \le 1 + \varepsilon_c \le \Lambda_0 < \infty$
follows from $A^{\mathrm{arith}} \preceq (1+\varepsilon_c)I$ and the transfer lemma. \qed
\end{proof}

\begin{corollary}[Perturbative universality]\label{cor:universality}
The gap $\lambda_{\min}(\GN)$ satisfies:
\begin{equation}\label{eq:universality}
  \lambda_{\min}(\GN)
  = \FTW + O\!\left(\left(\frac{\log\log c}{\log c}\right)^{1/2}\right)
  \quad\text{as }c\to\infty.
\end{equation}
The leading term $\FTW \approx 0.9597$ is \emph{universal}:
it depends only on the Airy kernel, not on the prime distribution.
The arithmetic correction is of order $o(1)$ and vanishes unconditionally.
\end{corollary}

\begin{proof}
Combine the lower bound \eqref{eq:main} with the upper bound
$\lambda_{\min}(\GN) \le \dAi + \varepsilon_c$. \qed
\end{proof}

%% =================================================================
\section{What the Proof Does Not Use}\label{sec:notuse}
%% =================================================================

We record precisely which structures are absent from the proof of
Theorem~\ref{thm:P13}, as this may be relevant for generalizations.

\begin{remark}[Logically absent structures]
The following appear in the companion papers (Papers I--XXI of this
programme) and in the broader PSWF literature,
but are \emph{not used} in the proof of Theorem~\ref{thm:P13}:
\begin{enumerate}[\rm (a)]
  \item The IIKS integrable-systems identity for the prolate kernel.
  \item The Riemann--Hilbert formulation and Deift--Zhou steepest descent.
  \item Widom's edge phase $\phi_c$ and the cubic Airy phase $\tfrac{1}{3}u^3$.
  \item Conjecture~3.3 of Paper~XII (convergence $Y_c \to Y_{\mathrm{Airy}}$
    in the RHP norm).
  \item The Generalized Riemann Hypothesis.
  \item Any conjecture on primes in short intervals
    stronger than the Montgomery--Vaughan mean square theorem.
\end{enumerate}
Items (a)--(d) are used in the companion papers to improve the
rate of coercivity from $c^{-1/2}$ to $c^{-1/3}$.
That improvement requires establishing
$\norm{K_c - \KAi}_{S_1} = O(c^{-1/3})$ (conditional on item~(d)),
but it does not affect the \emph{existence} of the gap established here.
\end{remark}

%% =================================================================
\appendix
%% =================================================================

\section{Proof of the Transfer Lemma}\label{app:transfer}

\begin{lemma}[Spectral transfer]\label{lem:transfer}
For all $c \ge c_1$ (explicit from PSWF Airy asymptotics):
\begin{equation}\label{eq:transfer}
  \lambda_{\min}(\GN)
  \ge \inf\spec(\DAi) - C_1 c^{-1/3}.
\end{equation}
\end{lemma}

\begin{proof}
By the Widom edge asymptotics \cite{Widom1964},
the rescaled PSWF Gram kernel converges in operator norm to the Airy kernel
on the window $\mathcal{W}$:
\begin{equation}
  \norm{\GN|_{\mathcal{W}} - \PAi A^{\mathrm{arith}} \PAi}_{\mathrm{op}}
  = O(c^{-1/3}).
\end{equation}
This is a consequence of the uniform Airy asymptotics for the PSWF
eigenvalues $\mu_k(c)$ in the transition zone:
$|\mu_k(c) - F_{\mathrm{TW}}((k - c^{2/3})/c^{1/3})| = O(c^{-1/3})$
uniformly for $k \in \mathcal{W}$ (see \cite{Widom1964,Bornemann2010}).
The claim follows by applying Weyl's inequality to the perturbation
$\GN|_{\mathcal{W}} - \PAi A^{\mathrm{arith}} \PAi$. \qed
\end{proof}

\section{Proof of the Arithmetic Perturbation Bound}\label{app:mv}

\begin{proof}[Proof of Lemma~\ref{lem:mv}]
The operator norm $\norm{\PAi\Ec\PAi}_{\mathrm{op}}$ is controlled by
the discrepancy of the prime counting function in the Airy window
$\mathcal{W} = \{N-K,\ldots,N-1\}$.

\emph{Step 1.} (Mode-counting in the Airy window).
Rescaling $k \mapsto t_k = (k - c^{2/3})/c^{1/3}$ for $k \in \mathcal{W}$,
the prime-indexed Airy nodes are
$p_j^{\mathrm{Ai}} := (p_j - c^{2/3})/c^{1/3}$
where $\{p_j\}$ are the primes in $[c^{2/3} - Kc^{1/3}, c^{2/3}]$.
The window has length $L \sim Kc^{1/3}$ around $x_0 = c^{2/3}$.

\emph{Step 2.} (Montgomery--Vaughan control).
By the Montgomery--Vaughan mean square theorem \cite[Theorem~7.3]{MV1974}:
\begin{equation}
  \sum_{j} \left(\pi(p_j^{\mathrm{Ai}} + h) - \pi(p_j^{\mathrm{Ai}}) - \frac{h}{\log p_j^{\mathrm{Ai}}}\right)^2
  \ll L \cdot \frac{\log\log x_0}{\log x_0}
  \qquad (h \ll L).
\end{equation}
This controls the $\ell^2$ norm of the fluctuation of the prime
density at the Airy scale.

\emph{Step 3.} (Operator norm lifting).
The operator $\Ec = A^{\mathrm{arith}} - I$ has kernel
$E_c(t_j, t_k) = A^{\mathrm{arith}}(t_j,t_k) - \delta_{jk}$
which is bounded in Hilbert--Schmidt norm by the $\ell^2$ discrepancy estimate.
The Airy-window Bernstein inequality (standard, see \cite{Young2001})
lifts the $\ell^2$ bound to an operator norm bound with the same rate.
This gives \eqref{eq:mv}. \qed
\end{proof}

%% =================================================================
\begin{thebibliography}{99}

\bibitem{TW1994}
C.~A.~Tracy and H.~Widom,
\textit{Level-spacing distributions and the Airy kernel},
Comm.\ Math.\ Phys.\ \textbf{159} (1994), 151--174.

\bibitem{Bornemann2010}
F.~Bornemann,
\textit{On the numerical evaluation of Fredholm determinants},
Math.\ Comp.\ \textbf{79} (2010), 871--915.

\bibitem{MV1974}
H.~L.~Montgomery and R.~C.~Vaughan,
\textit{Hilbert's inequality},
J.\ London Math.\ Soc.\ (2) \textbf{8} (1974), 73--82.

\bibitem{Kato1966}
T.~Kato,
\textit{Perturbation Theory for Linear Operators},
Springer, Berlin, 1966.

\bibitem{Widom1964}
H.~Widom,
\textit{Asymptotic behavior of the eigenvalues of certain integral operators},
Arch.\ Ration.\ Mech.\ Anal.\ \textbf{17} (1964), 215--229.

\bibitem{LandauWidom1980}
H.~J.~Landau and H.~Widom,
\textit{Eigenvalue distribution of time and frequency limiting},
J.\ Math.\ Anal.\ Appl.\ \textbf{77} (1980), 469--481.

\bibitem{Young2001}
R.~M.~Young,
\textit{An Introduction to Nonharmonic Fourier Series}, revised ed.,
Academic Press, San Diego, 2001.

\end{thebibliography}

\end{document}
